\section{Large girth outside a small set}\label{sec:half}

Recall that $q(G,P)$ is the probability that $\xi\cap E(G-P)$ is a forest.

\begin{theorem}\label{thm:half}
There are integers $D_*\ge7$ and $r_5$ with the following property. Let $G$ be a connected cubic
multigraph of rank $r\ge r_5$. Let $P$ be a nonempty set of at most $3\log r$ vertices such that
$G[P]$ is connected. If $G-P$ has girth at least $D_*$, then $q(G,P)<\frac25$.
\end{theorem}

\begin{idea}{outline of the proof}
Suppose that $q(G,P)\ge\frac25$. We find two disjoint parts of $G-P$ that are large, have small
boundaries and have large girth. Each of them has a forest proportion of at most $\frac35$, for
every choice of the edges outside it. This is a contradiction, because $(\frac35)^{2}<\frac25$. To
find the parts we use a packing bound for short cycles (Section~\ref{sec:packing}), a second
moment and a chain of nested separators (Section~\ref{sec:chain}).
\end{idea}

\subsection{Pieces with a small boundary}\label{sec:affine}

Let $N$ be a pole, and let $T$ be a set of ports of even size. A \emph{$T$-join} of $N$ is a set
of edges in which the vertices of $T$ have odd degree and all other vertices have even degree.
There are $2^{\beta(N)}$ of them. The \emph{forest proportion} of $T$ is the proportion of the
$T$-joins that are forests. For $T=\emptyset$ the $T$-joins are the even subgraphs, only the empty
one is a forest, and the forest proportion is $2^{-\beta(N)}$. Put
\[
  R_{*}=R_{1/20},\qquad B=2R_{*}+2^{100}+1,
\]
with $R_{1/20}$ as in Proposition~\ref{prop:input}. Then $\Theta_{1/20}(k)+k\le B(k+1)^{2}$ for all
$k\ge1$.

\begin{lemma}[A marker in a pole]\label{lem:marker}
Let $N$ be a pole with $k\ge2$ ports, and let $o$ be a port. Call the set of the vertices at
distance $j$ from $o$ the \emph{layer} $j$. Suppose that some set $T$ of ports has forest
proportion at least $\frac35$. Then each layer has at most $B(k+1)^{2}$ vertices.
\end{lemma}

\begin{proof}
Take a breadth-first tree of $N$ from $o$, and let $P$ be the union of its paths from $o$ to the
other ports. Then $N[P]$ is connected, $P$ contains all vertices of degree two, and $P$ has at most
$k$ vertices in each layer, because a path of the tree from $o$ has one vertex in each layer that
it meets. The tree paths inside $P$ contain a $T$-join $F_0$. The map $J\mapsto J+F_0$ is a
bijection from the $T$-joins of $N$ to the even subgraphs of $N$, and $J+F_0$ has the same edges in
$N-P$ as $J$. So if $J$ is a forest, then $J+F_0$ is a forest in $N-P$. Hence
$q(N,P)\ge\frac35>\half+\frac1{20}$. The distance from $o$ is a level function. By
Lemma~\ref{lem:width} with $\delta=\frac1{20}$, each layer has fewer than $\Theta_{1/20}(k)$
vertices outside $P$, and it has at most $k$ vertices in $P$.
\end{proof}

\begin{lemma}[Girth and rank of a pole]\label{lem:moore}
Let $N$ be a pole of rank $h\ge2$ with $k$ ports and girth at least $D$. Put
$t=\lceil\log(8h)\rceil$. Then
\[
  D\;\le\;4t\Bigl(1+\frac{k}{3(h-1)}\Bigr).
\]
\end{lemma}

\begin{proof}
Replace every vertex of degree two and its two edges by one edge, and give each edge of the result
the length of the path of $N$ that it replaces. The result $N_0$ is a connected cubic multigraph
with $n_0=2h-2$ vertices and $m_0=3h-3$ edges, and the lengths of its edges add up to $m_0+k$. A
cycle of $N_0$ comes from a cycle of $N$ of the same length.

There are $2m_0\,2^{t-1}$ non-backtracking walks with $t$ edges in $N_0$. Choose one of them
uniformly. Then each of its steps is uniform over the $2m_0$ directed edges. So the expected length
of the walk is $t(1+k/m_0)$. At least half of the walks have length at most twice this value. Their
number is at least
\[
  m_0\,2^{t-1}\;\ge\;12h(h-1)\;>\;n_0^{2}.
\]
If one of these short walks visits a vertex twice, it contains a cycle that is not longer than the
walk. Otherwise all of them are paths, and two of them have the same first vertex and the same last
vertex. Their union contains a cycle of length at most the sum of their lengths. In both cases $N$
has a cycle of length at most $4t(1+k/m_0)$.
\end{proof}

\begin{lemma}[Forest proportions at large girth]\label{lem:affine}
There are absolute constants $D_{*}\ge7$ and $C_0$ with the following property. Let $N$ be a pole
of rank $h$ with $k$ ports and girth at least $D$. Suppose that
\begin{enumerate}
\item[\textup{(a)}] $h\ge256$, $k\le4D$ and $D\ge D_{*}$, or
\item[\textup{(b)}] $h\ge\Lambda\ge4$, $k\le2\Lambda$ and $D\ge C_0\log(\Lambda+2)$ for some
number $\Lambda$.
\end{enumerate}
Then every set of ports of even size has forest proportion at most $\frac35$.
\end{lemma}

\begin{proof}
If $k\le1$, the only set is the empty one, with forest proportion $2^{-h}$. Let $k\ge2$, and put
$t=\lceil\log(8h)\rceil\le\log h+4$.

\step{The rank is exponential in the girth.} In case~(a), $h-1>\frac{32}3t$: this holds at
$h=256$, where $t=11$, and the ratio $(h-1)/(\log h+4)$ increases with $h$. So
Lemma~\ref{lem:moore} gives $D\le4t+16tD/(3(h-1))<4t+D/2$. In case~(b),
$k/(3(h-1))\le2h/(3(h-1))<1$, and Lemma~\ref{lem:moore} gives $D<8t$. In both cases
$D<8t\le8\log h+32$, so
\begin{equation}\label{eq:rankgirth}
  h\;\ge\;2^{D/8-4},\qquad\text{and }N\text{ has }n\ge h\text{ vertices}.
\end{equation}

\step{A ball without ports.} Put $s=\lfloor D/16\rfloor$. A ball of radius $s$ has at most
$3\cdot2^{s}$ vertices. So at most $3k\,2^{s}$ vertices are within distance $s$ of a port. In
case~(a) this is at most $12D\,2^{D/16}$, and in case~(b) at most $6\Lambda\,2^{D/16}$. Choose
$D_{*}$ and $C_0$ so large that this number is less than $2^{D/8-4}\le n$ in each case. Then some
vertex $w$ has distance more than $s$ from every port. All vertices within distance $s$ of $w$
have degree three, and $2s<D$. So the ball of radius $s$ around $w$ is a tree in which every
vertex at distance less than $s$ from $w$ has degree three, and it has at least
$2^{s}\ge2^{D/16-1}$ vertices. Fix a port $o$. The distances from $o$ to two vertices of this ball
differ by at most $2s$. So the ball meets at most $2s+1\le D$ layers with respect to $o$, and some
layer has at least $2^{D/16-1}/D$ vertices.

\step{Conclusion.} Suppose that some set of ports has forest proportion more than $\frac35$. By
Lemma~\ref{lem:marker} every layer has at most $B(k+1)^{2}$ vertices. This is at most
$B(4D+1)^{2}$ in case~(a) and at most $B(2\Lambda+1)^{2}$ in case~(b). In case~(a) choose $D_{*}$ so
large that $2^{D/16-1}/D>B(4D+1)^{2}$ for $D\ge D_{*}$. In case~(b) the function $2^{D/16}/D$
increases for $D\ge32$, and at $D=C_0\log(\Lambda+2)$ its ratio to $(\Lambda+2)^{2}$ is
$(\Lambda+2)^{C_0/16-2}/(C_0\log(\Lambda+2))$. This tends to infinity with $C_0$, uniformly in
$\Lambda\ge4$. So a large $C_0$ gives $2^{D/16-1}/D>B(2\Lambda+1)^{2}$. In both cases we have a
contradiction.
\end{proof}

\paragraph{Products of pieces.}

\begin{lemma}[Product]\label{lem:product}
Let $G$ be a cubic multigraph, let $P$ be a set of vertices, and let
$U_1,\ldots,U_m\subseteq V(G)\sm P$ be pairwise disjoint sets such that each $G[U_i]$ is a pole.
Suppose every even set of ports in each pole has forest proportion at most $3/5$.
Then $q(G,P)\le(3/5)^m$.
\end{lemma}

\begin{proof}
Let $E'$ consist of the edges that lie in none of the induced poles.
Fix a restriction $x\subseteq E'$ that extends to an even subgraph of $G$.
For each $i$, let $T_i$ be the vertices of $U_i$ incident with an odd number of edges of $x$.
The parity equations say that $\xi\cap E(G[U_i])$ is a $T_i$-join.
These equations involve disjoint internal edge sets, so the joins are independent and uniform,
conditional on $\xi\cap E'=x$.
A vertex of degree three in a pole has no incident edge in $E'$; thus each $T_i$ is an even set of ports.
If $\xi\cap E(G-P)$ is a forest, its restriction to every pole is a forest.
The conditional probability is therefore at most $(3/5)^m$. Average over $x$.
Edges between the poles belong to $E'$, so adjacent pieces cause no exception.
\end{proof}

\begin{lemma}[The core of a piece]\label{lem:core}
Let $G$ be a cubic multigraph, and let $X$ be a set of vertices such that $G[X]$ is connected, has
no loop, and has rank $h\ge2$. Then there is a set $U\subseteq X$ such that $G[U]$ is a pole of rank $h$ with at
most $|\delta X|$ ports, and the girth of $G[U]$ is at least the girth of $G[X]$.
\end{lemma}

\begin{proof}
Delete vertices of degree one from $G[X]$ as long as there are any. The result $G[U]$ is connected,
has the same cycles as $G[X]$ and all its degrees are two or three. When a vertex of degree one is
deleted, two edges that left the set no longer do, and one edge inside the set now leaves it. So
$|\delta U|\le|\delta X|$. A vertex of degree two in $G[U]$ has exactly one edge end outside
$G[U]$, and a vertex of degree three has none. So $G[U]$ has $|\delta U|$ ports.
\end{proof}

\begin{proposition}[Pieces with small boundaries]\label{prop:blocks}
Let $G$ be a connected cubic multigraph and let $P$ be a set of vertices. Let
$X_1,\ldots,X_m\subseteq V(G)\sm P$ be pairwise disjoint, nonempty and proper sets of vertices. Suppose that one of
the following holds for every $i$.
\begin{enumerate}
\item[\textup{(a)}] $G[X_i]$ is connected of rank at least $256$ and girth at least $D\ge D_{*}$,
and $|\delta X_i|\le4D$.
\item[\textup{(b)}] $|X_i|\ge8\Lambda^{2}$, $|\delta X_i|\le2\Lambda$, and $G[X_i]$ has girth at
least $D$, where $\Lambda\ge4$ and $D\ge C_0\log(\Lambda+2)$.
\end{enumerate}
Then $q(G,P)\le(3/5)^m$.
\end{proposition}

\begin{proof}
It suffices to find sets $U_i\subseteq X_i$ as in Lemma~\ref{lem:product}. In case~(a),
Lemma~\ref{lem:core} gives a pole $G[U_i]$ of rank at least $256$ with at most $4D$ ports and girth
at least $D$, and Lemma~\ref{lem:affine}(a) applies.

In case~(b) write $n=|X_i|$, $d=|\delta X_i|$ and $c=\comp(G[X_i])$. Every component of $G[X_i]$
has an edge that leaves $X_i$, so $1\le c\le d\le2\Lambda$. Count the edge ends:
$3n=2|E(G[X_i])|+d$. So $\beta(G[X_i])=(n-d)/2+c$, and some component $Y$ of $G[X_i]$ has rank at
least
\[
  \frac{n-d}{2c}+1\;\ge\;\frac{8\Lambda^{2}-2\Lambda}{4\Lambda}+1\;>\;\Lambda .
\]
Its boundary is a part of $\delta X_i$. Lemma~\ref{lem:core} for $Y$ gives a pole $G[U_i]$ of rank
at least $\Lambda$ with at most $2\Lambda$ ports and girth at least $D$, and
Lemma~\ref{lem:affine}(b) applies.
\end{proof}

\subsection{Packing of short cycles}\label{sec:packing}

\begin{idea}{the purpose of the packing bound}
The second-moment argument requires a graph of girth about $\log r$. So we delete a maximal
family of disjoint short cycles, and we need a bound for the size of this family.
Lemma~\ref{lem:smallcuts} gives such a bound only when the proportion of forests exceeds one half
by a fixed margin. Here we prove the bound with Proposition~\ref{prop:blocks}. Many disjoint
cyclic sets with small boundaries have independent local events, or many of them lie in one
string (Lemma~\ref{lem:strings}), and a long string contains two pieces to which
Proposition~\ref{prop:blocks} applies.
\end{idea}

Throughout, $h_*=256$.

\paragraph{Vertices that share a parity.}

Let $Q$ be a connected multigraph without loops. For an edge $e$, the map
$\xi\mapsto[e\in\xi]$ is a linear functional on $\Z(Q)$. For a vertex $v$ let $L_v$ be the space
of functionals that is spanned by the edges at $v$. For a set $A$ of vertices put
$U_A=\sum_{a\in A}L_a$. We use one standard fact: if a set $F$ of edges meets every even subgraph of
$Q$ in an even number of edges, then $F=\delta Y$ for some set $Y$ of vertices. (Fix a spanning
tree and a root. Let $Y$ be the set of the vertices whose tree path from the root has an odd number
of edges in $F$. Then $F$ and $\delta Y$ contain the same tree edges, and the fundamental cycles
show that they contain the same edges outside the tree.)

\begin{lemma}[Strings]\label{lem:strings}
Let $Q$ be a connected multigraph without loops, and let $A$ be a nonempty proper set of its
vertices. Put $b=|\delta A|$. Call a vertex $v\notin A$ \emph{relevant} if $L_v\cap U_A\ne0$. There
are a set of at most $2b$ relevant vertices, the \emph{exceptions}, and at most $2b$ sequences of
relevant vertices, the \emph{strings}, with these properties. Every relevant vertex that is not an
exception is in a string. If $v_1,\dots,v_m$ is a string and $1\le i<j\le m$ with $j\ge i+2$, there
is a set $M(i,j)$ of vertices such that
\begin{enumerate}
\item[\textup{(i)}] $M(i,j)$ is disjoint from $A\cup\{v_i,v_j\}$, the graph $Q[M(i,j)]$ is
connected, and $v_l\in M(i,j)$ for $i<l<j$;
\item[\textup{(ii)}] every edge that leaves $M(i,j)$ ends in $v_i$ or in $v_j$;
\item[\textup{(iii)}] $M(i,j)\cap M(i',j')=\emptyset$ if $j\le i'$.
\end{enumerate}
\end{lemma}

\begin{proof}
Call the edges of $\delta A$ \emph{terminals}. The end of a terminal outside $A$ is its
\emph{foot}. Every component of $Q-A$ contains a foot, because $Q$ is connected.

\step{Step 1.} Let $v\notin A$ be relevant and not a foot, and let $K$ be the component of $Q-A$
that contains $v$. We show that at least two components of $K-v$ contain feet. Let
$0\ne a\in L_v\cap U_A$. Then $a$ is the sum of the functionals of a set $F$ of edges at $v$, and
also of a set $F'$ of edges with an end in $A$. No edge is in both sets, because $v$ is not a foot.
The set $F\cup F'$ meets every even subgraph in an even number of edges, so $F\cup F'=\delta Y$ for
some $Y$. We may assume that $v\notin Y$. No edge of $K-v$ is in $F\cup F'$. So each component of
$K-v$ lies inside $Y$ or outside $Y$, and $F$ is the set of all edges from $v$ to the components
inside $Y$. For a component $J$ of $K-v$, the cut $\delta V(J)$ consists of the edges from $J$ to
$v$ and the terminals with a foot in $J$. So on $\Z(Q)$ the sum of the edges from $J$ to $v$ equals
the sum of the terminals with a foot in $J$. Hence $a$ is the sum of the terminals with a foot in
$Y$. The sum of all terminals with a foot in $K$ is zero on $\Z(Q)$, because these terminals form
the cut $\delta V(K)$. As $a\ne0$, some component of $K-v$ inside $Y$ contains a foot, and some
component outside $Y$ contains a foot.

\step{Step 2.} In each component $K$ of $Q-A$ fix a spanning tree. For each terminal add a new leaf
at its foot. Let $T_K$ be the union of the paths of this tree between the new leaves, and let $T$
be the union of the trees $T_K$. Every leaf of $T$ is a new leaf, so $T$ has at most $b$ leaves and
at most $b$ vertices of degree at least three. The exceptions are the relevant vertices that are
feet or have degree at least three in $T$. There are at most $2b$ of them.

Let $v$ be a relevant vertex that is not an exception. By Step~1 two components of $K-v$ contain
feet. The path of the tree between the two new leaves at these feet passes through $v$. So $v$ is
a vertex of $T$, and its degree in $T$ is two. A \emph{thread} is a maximal path of $T$ whose inner
vertices have degree two in $T$. There are at most $2b$ threads, and every vertex of degree two of
$T$ is an inner vertex of exactly one thread. The strings are the sequences of the relevant
vertices that are not exceptions on one thread, in the order of the thread.

\step{Step 3.} Let $v$ be a vertex of a string, in the component $K$. The forest $T-v$ has exactly
two components that meet $K$, one on each side of $v$ along its thread. Their vertex sets in $K$ are
the \emph{near side} of $v$, towards the start of the thread, and the \emph{far side} of $v$. Each
side is connected in $K-v$, and all feet in $K$ lie in the two sides. By Step~1 the two sides lie in
different components of $K-v$.

Now let $v_i$ and $v_j$ be as in the lemma. Let $\Pi$ be the part of their thread strictly between
them. It contains $v_{i+1}$, it is connected, and it contains no foot, because a new leaf at a
vertex of $\Pi$ would give that vertex degree three in $T$. Let $M(i,j)$ be the vertex set of the
component of $K-v_i-v_j$ that contains $\Pi$. This gives~(i).

For~(ii), suppose that $M(i,j)$ contains a foot $w$. Then $w$ is a vertex of $T$ outside $\Pi$. So
$w$ lies in the near side of $v_i$ or in the far side of $v_j$; say in the near side of $v_i$. The
path $\Pi$ lies in the far side of $v_i$. But $w$ and $\Pi$ are joined in $K-v_i$ through
$M(i,j)$, and the two sides of $v_i$ lie in different components of $K-v_i$. This is a
contradiction. So $M(i,j)$ contains no foot. An edge that leaves $M(i,j)$ ends in $A$, in $v_i$ or
in $v_j$, and it does not end in $A$.

For~(iii), let $j\le i'$. The set $M(i,j)$ is connected in $K-v_j$ and contains $\Pi$, which lies
in the near side of $v_j$. So $M(i,j)$ lies in the component of $K-v_j$ that contains the near side
of $v_j$. We show that $M(i',j')$ lies in the component of $K-v_j$ that contains the far side of
$v_j$; then the two sets are disjoint. The set $M(i',j')$ is connected in $K-v_{i'}$ and contains
the part $\Pi'$ of the thread between $v_{i'}$ and $v_{j'}$, which lies in the far side of
$v_{i'}$. If $j=i'$, this is the claim. If $j<i'$, then $v_j$ lies in the near side of $v_{i'}$, so
$v_j\notin M(i',j')$. Hence $M(i',j')$ is connected in $K-v_j$, and it contains $\Pi'$, which lies
in the far side of $v_j$.
\end{proof}

\paragraph{A packing bound.}

\begin{proposition}[Packing]\label{prop:packing}
Let $G$ be a connected cubic multigraph. Let $P$ be empty or induce a connected subgraph,
and let $S$ consist of $s$ vertices outside $P$.
Suppose $G-(P\cup S)$ has girth at least $D\ge D_*$ and $q(G,P\cup S)\ge\frac25$.
Let $2\le d\le2D$, and let $X_1,\ldots,X_n$ be pairwise disjoint sets outside $P\cup S$,
each inducing a connected graph with a cycle and satisfying $|\delta X_i|\le d$.
Then $n\le1551\,d(s+1)2^{d}$.
\end{proposition}

\begin{proof}
Contract each $X_i$ to a marked vertex $x_i$, contract each component of
$G-\bigcup_iX_i$ to a vertex, and delete the loops. The resulting graph $Q$ is connected.
Projection maps the uniform distribution on $\Z(G)$ to the uniform distribution on $\Z(Q)$,
as in Lemma~\ref{lem:loc}(b). Conditional on the projected edges, the restrictions to the
induced graphs $G[X_i]$ are independent affine joins.
Let $I_i$ be the event that the projection uses no edge at $x_i$.
On $I_i$ the restriction to $G[X_i]$ is an even subgraph, and its forest probability is at most
$1/2$, since $G[X_i]$ contains a cycle. Hence
\begin{equation}\label{eq:packprod}
 \frac25\le q(G,P\cup S)\le
 \EE\prod_{i=1}^n\left(1-\frac{\mathbf1_{I_i}}2\right).
\end{equation}

Write $L_i=L_{x_i}$ for the local parity space in $Q$.
The sum of its incident edge functionals is zero, so $\dim L_i\le d-1$ and
$\PP(I_i)=2^{-\dim L_i}\ge2^{1-d}$.
If $z_0$ local spaces are zero, then $\frac25\le2^{-z_0}$.
Thus at most one local space is zero.
Keep its vertex in $Q$, but omit this mark from the selection below.
If no mark remains, the desired bound is immediate.

Choose a maximal direct-sum family $A$ of nonzero local spaces.
Put $t=|A|\ge1$ and $b=|\delta A|\le dt$.
The events $I_a$ for $a\in A$ are independent, so
\[
 \frac25\le\prod_{a\in A}\left(1-\frac{\PP(I_a)}2\right)
 \le(1-2^{-d})^t<e^{-t2^{-d}}.
\]
In particular, $t<2^d$.
If all remaining marks lie in $A$, the bound follows.
Otherwise $A$ is a nonempty proper vertex set of $Q$.
Every remaining mark outside $A$ is relevant in Lemma~\ref{lem:strings}, by maximality.
There are at most $2b$ exceptions and at most $2b$ strings.

Suppose a string has at least $(h_*+2)(s+3)$ marks.
Split these marks into $s+3$ consecutive groups of $h_*+2$ marks.
For each group take the region between its first and last marks from Lemma~\ref{lem:strings}.
Lift this region to the union of the sets contracted to its vertices.
The lifted regions are connected and pairwise disjoint.
Each contains $h_*$ disjoint cyclic sets $X_i$, so has rank at least $h_*$.
Its boundary has at most $2d\le4D$ edges.
At most $s$ regions meet $S$, and at most one meets $P$, since a nonempty connected $P$
lies in one contracted set. Thus two regions remain inside $G-(P\cup S)$.
Proposition~\ref{prop:blocks}(a) gives $q(G,P\cup S)\le\frac9{25}<\frac25$, a contradiction.

Count the omitted mark, the selected marks, the exceptions and the strings:
\[
 n\le1+t+2b+2b(h_*+2)(s+3).
\]
Here $b\le dt$, $t<2^{d}$, $s+3\le3(s+1)$ and $h_*+2=258$. So
$n\le2^{d}+2d\,2^{d}+1548\,d(s+1)2^{d}\le1551\,d(s+1)2^{d}$.
\end{proof}

\paragraph{Raising the girth.}

\begin{theorem}[Raising the girth]\label{thm:raise}
Let $G$ be a connected cubic multigraph and let $P$ be empty or induce a connected subgraph.
Suppose $q(G,P)\ge\frac25$ and $G-P$ has girth at least $D_*$.
Put $D_j=2^jD_*$ for $j\ge0$.
There are nested sets $S_0\subseteq S_1\subseteq\cdots$ outside $P$ such that
\begin{enumerate}
\item[\textup{(i)}] $S_j$ is the union of a maximal family of pairwise disjoint cycles of length
at most $D_j$ in $G-P$; hence $G-(P\cup S_j)$ has girth greater than $D_j$;
\item[\textup{(ii)}] $\log(|S_j|+1)\le2D_j+6\log^{2}D_j$.
\end{enumerate}
\end{theorem}

\begin{proof}
Put $S_{-1}=\emptyset$.
Having chosen $S_{j-1}$, take a maximal disjoint family $\mathcal Q_j$ of cycles of length
at most $b=D_j$ in $G-(P\cup S_{j-1})$, and add its vertices to $S_{j-1}$ to form $S_j$.
The exterior girth is at least $a=D_*$ for $j=0$, and at least $a=D_{j-1}+1$ otherwise;
thus $b\le2a$.
Each cycle $C$ has $|\delta V(C)|\le|C|\le b$.
Also $q(G,P\cup S_{j-1})\ge q(G,P)\ge\frac25$.
Proposition~\ref{prop:packing} gives
\[
 |\mathcal Q_j|\le1551\,b(|S_{j-1}|+1)2^b,
 \qquad |S_j|+1\le1552\,b^22^b(|S_{j-1}|+1).
\]
The union of all families chosen so far is maximal at scale $D_j$, which proves~(i).
For~(ii), take logarithms in the recursion. As $\log1552<11$, the step $i$ adds at most
$D_i+2\log D_i+11$. There are $j+1\le\log D_j$ steps, and $\sum_{i\le j}D_i<2D_j$. So
\[
 \log(|S_j|+1)\le2D_j+\log D_j\,\bigl(2\log D_j+11\bigr)\le2D_j+6\log^{2}D_j,
\]
because $\log D_j\ge\log7>\frac{11}4$.
\end{proof}

\subsection{Links, a second moment and a chain}\label{sec:chain}

\paragraph{Components between two connected sets.}
Lemma~\ref{lem:smallcuts} bounds the number of disjoint cyclic regions with small boundaries, and
it needs an excess above $\half$. For components that join the same two connected sets there is a
bound without an excess.

\begin{lemma}[Parallel links]\label{lem:parallel}
Let $\Gamma$ be a connected multigraph, let $A$ and $B$ be disjoint regions, and let $P$ be a set of
vertices with $q(\Gamma,P)\ge\frac25$. Let $t\ge0$ be an integer. Then at most $5\cdot2^{t}$
components $Y$ of $\Gamma-(A\cup B)$ have all of the following properties: $Y$ has an edge to $A$
and an edge to $B$; $Y$ avoids $P$; $Y$ contains a cycle; $|\delta Y|\le t$.
\end{lemma}

\begin{proof}
Let $Y_0,Y_1,\dots,Y_n$ be such components. For $1\le i<j\le n$ the graph $\Gamma-Y_i-Y_j$ is
connected. Indeed, it contains the connected set $A\cup B\cup Y_0$, and every other component of
$\Gamma-(A\cup B)$ is connected and has an edge to $A\cup B$. No edge joins $Y_i$ and $Y_j$. So the
regions $Y_i$ and $Y_j$ have exactly one link, and by Lemma~\ref{lem:link}
\[
  \PP(\cO_{Y_i}\cap\cO_{Y_j})=\PP(\cO_{Y_i})\,\PP(\cO_{Y_j}).
\]
Let $E_i$ be the event that $\cO_{Y_i}$ holds and $\xi\cap E(\Gamma[Y_i])$ is not empty. By the last
part of Lemma~\ref{lem:link}, the events $E_i$ and $E_j$ are independent as well, and
$\PP(E_i)\ge2^{-t}\cdot\half$. If $n=0$, the bound holds. Otherwise let $N$ be the number of the
events $E_1,\dots,E_n$ that occur, and put $\nu=\EE N\ge n\,2^{-t-1}>0$. Then
$\EE N^{2}\le\nu^{2}+\nu$, and the Cauchy--Schwarz inequality gives
\[
  \PP(N=0)\;\le\;1-\frac{\nu^{2}}{\EE N^{2}}\;\le\;\frac1{1+\nu}.
\]
If $\xi\cap E(\Gamma-P)$ is a forest, then no $E_i$ occurs, because $Y_i$ avoids $P$ and a nonempty
even subgraph is not a forest. So $\frac25\le q(\Gamma,P)\le1/(1+\nu)$. Hence $\nu\le\frac32$ and
$n\le3\cdot2^{t}$. With $Y_0$ there are at most $3\cdot2^{t}+1$ components.
\end{proof}

\begin{lemma}[Few links]\label{lem:linkshalf}
Let $G$ be a connected cubic multigraph, and let $P$ be a nonempty set of vertices such that $G[P]$
is connected and $q(G,P)\ge\frac25$. Let $S$ be the union of the vertex sets of a maximal family of
pairwise disjoint cycles of length at most $D$ in $G-P$, and put $J=G-(P\cup S)$. Let
$8\le D\le L$, and let $t$ be an integer with $1\le t\le D/4$. Then any two disjoint cycles $C,C'$
of $J$ of length at most $L$ satisfy
\[
  \ell\bigl(V(C),V(C')\bigr)\;\le\;1+\frac{2L}{t+1}+5\cdot2^{t}+289\Bigl(\frac LD\Bigr)^{2}.
\]
\end{lemma}

\begin{proof}
By the maximality of the family, $J$ is a simple graph of girth greater than $D$. So we are in the
setting of Section~\ref{sec:moderate}, with $\Gamma=G$ and $U=P\cup S$. Put
$U'=V(C)\cup V(C')$. Every vertex of $U'$ lies on exactly one edge outside its own cycle, so at
most $2L$ edges leave $U'$. At most one component of $G-U'$ meets $P$, because $G[P]$ is connected
and disjoint from $U'$. At most $2L/(t+1)$ components have more than $t$ edges to $U'$. Of the
other components that are joined to both cycles, at most $5\cdot2^{t}$ contain a cycle, by
Lemma~\ref{lem:parallel} with $A=V(C)$ and $B=V(C')$. The remaining components that are joined to
both cycles induce trees, avoid $P$ and have at most $t$ edges to $U'$. They also avoid $S$:
each cycle of the family is disjoint from $U'$, so it lies in one component of $G-U'$, and this
component contains a cycle. By Lemma~\ref{lem:treelinks}, at most $289(L/D)^{2}$ links are such
components or edges.
\end{proof}

\paragraph{Pairs with exactly two links.}
Let $G$, $P$, $S$, $J$, $D$ and $L$ be as in Lemma~\ref{lem:linkshalf}. As in its proof we are in
the setting of Section~\ref{sec:moderate}. So $s=\lfloor(D-1)/2\rfloor$, the set $\mathcal{C}$
consists of the cycles of $J$ of length at most $L$, and the numbers $a_C$, $\mu$ and $W_2$ and the
random variables $X_C$ and $Z$ are defined there.

\begin{lemma}[Two links]\label{lem:twolinks}
Suppose that $\ell(V(C),V(C'))\le b$ for all disjoint $C,C'\in\mathcal{C}$, and put
$\eta=2^{b}L^{3}2^{-s}$. Then
\[
  W_2\;\ge\;\tfrac23\,\mu^{2}-(1+\eta)\,\mu .
\]
\end{lemma}

\begin{proof}
By Lemma~\ref{lem:moment}, $\EE Z^{2}\le\mu^{2}+W_2+(1+\eta)\mu$. If $\xi\cap E(G-P)$ is a forest,
then $\xi$ contains no cycle of $J$, so $Z=0$. Hence $\PP(Z\ne0)\le\frac35$. The Cauchy--Schwarz
inequality $\mu^{2}\le\EE Z^{2}\cdot\PP(Z\ne0)$ gives $\EE Z^{2}\ge\frac53\mu^{2}$.
\end{proof}

\begin{idea}{from pairs with two links to a chain}
By Lemma~\ref{lem:twolinks} there are many pairs of cycles with exactly two links. Let $C$ and
$C'$ be such a pair. One of the two links avoids $p$, and then $C'$ separates a neighbour $u$ of
$C$ from $p$. Many disjoint cycles $C'$ have this property for the same $C$ and $u$. They are
nested (Figure~\ref{fig:chain}), and the parts between them are the pieces that we need. The
argument requires that $C'$ lies in the component of $p$ in $G-V(C)$. By symmetry, at least half
of the weight of the pairs with two links has this property.
\end{idea}

\paragraph{Separators are nested.}

\begin{lemma}[Nested separators]\label{lem:nested}
Let $u$ and $p$ be two vertices of a connected graph $H$. Let $C_1$ and $C_2$ be disjoint connected
sets of vertices that avoid $u$ and $p$, such that $u$ and $p$ lie in different components of
$H-C_i$ for $i=1,2$. Let $U_i$ be the vertex set of the component of $u$ in $H-C_i$. Then exactly
one of the following holds: $C_1\subseteq U_2$ and $U_1\subseteq U_2$; or $C_2\subseteq U_1$ and
$U_2\subseteq U_1$.
\end{lemma}

\begin{proof}
Take a path from $u$ to $p$, and let $z$ be its last vertex before the first vertex in $C_1$. The
part of the path up to $z$ lies in $U_1$. In the same way there is a path from $p$ to a vertex next
to $C_1$ that avoids $C_1$. The set $C_2$ is connected and disjoint from $C_1$, so it lies in one
component of $H-C_1$. If this component is neither $U_1$ nor the component of $p$, then the two
partial paths and a path inside $C_1$ form a walk from $u$ to $p$ that avoids $C_2$. This is
impossible. So $C_2$ lies in $U_1$ or in the component of $p$ in $H-C_1$, and the same holds with
the indices exchanged.

Suppose that $C_2$ lies in the component of $p$ in $H-C_1$. Then $C_2$ avoids $U_1$. So every
vertex of $U_1$ is joined to $u$ outside $C_2$, and so is every vertex of $C_1$, through $z$ and
the connected set $C_1$. Hence $U_1\cup C_1\subseteq U_2$. If $C_2$ lies in $U_1$ and $C_1$ lies in
$U_2$, then the path from $p$ to a vertex next to $C_1$, a path inside $C_1$, and a path inside
$U_2$ from $C_1$ to $u$ give a walk from $p$ to $u$ outside $C_2$. This is impossible. So if $C_2$
lies in $U_1$, then $C_1$ lies in the component of $p$ in $H-C_2$, and
$U_2\cup C_2\subseteq U_1$ by the first case with the indices exchanged. The two alternatives
exclude each other, because $C_1\subseteq U_2$ and $U_2\subseteq U_1$ would give
$C_1\subseteq U_1$.
\end{proof}

So a family of pairwise disjoint connected sets that separate $u$ from $p$ is linearly ordered:
write $C_1<C_2$ in the first case of the lemma. The order is transitive, because it is the order
of the sets $U_i$ by inclusion.

\begin{lemma}[Blocks]\label{lem:blocks}
In the situation of Lemma~\ref{lem:nested}, let $C_1<C_2$ and put
$X=U_2\sm(U_1\cup C_1)$. Then every edge of $H$ that leaves $X$ ends in $C_1$ or in $C_2$. If
$C_1<C<C_2$ for a third set $C$ of the family, then $C\subseteq X$. If $C_1<C_2\le C_3<C_4$, then
the sets $U_2\sm(U_1\cup C_1)$ and $U_4\sm(U_3\cup C_3)$ are disjoint.
\end{lemma}

\begin{proof}
An edge that leaves $U_2$ ends in $C_2$, because $U_2$ is a component of $H-C_2$. No edge joins
$X$ to $U_1$: its end in $X$ would lie in the component $U_1$ of $H-C_1$. So an edge from $X$ to
$U_1\cup C_1$ ends in $C_1$. For the second statement,
$C\subseteq U_2$, and $C$ avoids $U_1\cup C_1$ because $U_1\cup C_1\subseteq U_C$, the component of
$u$ in $H-C$. For the third, $U_2\subseteq U_3$.
\end{proof}

\paragraph{The chain.}

\begin{lemma}[Chain]\label{lem:chain}
In the situation of Lemma~\ref{lem:twolinks}, suppose that
\[
  \mu>30,\qquad\eta\le\frac12,\qquad L^{2}2^{-s}\le\frac12 .
\]
Fix a vertex $p\in P$. Then there are a cycle $C\in\mathcal C$, a vertex $u\notin V(C)$ that has a
neighbour on $C$, and pairwise disjoint cycles $C_1,\dots,C_T\in\mathcal C$ with
\[
  T\;\ge\;\frac{8\cdot2^{s}}{L^{3}},
\]
such that each $C_i$ is disjoint from $C$ and avoids $u$, and $u$ and $p$ lie in the same component
of $G-V(C)$ and in different components of $G-V(C)-V(C_i)$.
\end{lemma}

\begin{proof}
For disjoint $C,C'\in\mathcal C$ say that $C'$ is \emph{far behind} $C$ if $C'$ does not lie in the
component of $p$ in $G-V(C)$. This relation is antisymmetric. Let $C'$ be far behind $C$. A path
from $p$ to its first vertex on $C$ stays in the component of $p$ in $G-V(C)$ before it ends. So it
avoids $C'$. As $C$ is connected and disjoint from $C'$, all of $C$ lies in the component of $p$ in
$G-V(C')$. So $C$ is not far behind $C'$. Hence the ordered pairs $(C,C')$ with $C'$ far behind $C$
have total weight at most $W_2/2$ when restricted to the exactly-two-link pairs.
Indeed, the two-link relation is symmetric and opposite orientations have the same weight;
at most one orientation is far behind.

Fix $C$. Every vertex of $C$ has exactly one edge end outside $C$, so at most $L$ vertices outside
$C$ have a neighbour on $C$. A cycle $C'$ that is joined to $C$ by an edge passes through one of
them. By~\eqref{eq:load} these cycles have total weight at most $L^{2}2^{-s}\le\half$. So the
ordered pairs of disjoint cycles that are joined by an edge have total weight at most $\mu/2$.

Call an ordered pair $(C,C')$ \emph{good} if the cycles are disjoint, have exactly two links, are
not joined by an edge, and $C'$ is not far behind $C$. By Lemma~\ref{lem:twolinks} and the two
bounds above, the good pairs have total weight at least $W_2/2-\mu/2$, hence at least
\[
  \frac13\mu^2-\left(1+\frac\eta2\right)\mu
  \;>\;\left(10-\frac54\right)\mu\;>\;8\mu .
\]
So for some $C\in\mathcal C$ the cycles $C'$ with $(C,C')$ good have total weight more than $8$.
Fix this $C$, and put $H=G-V(C)$.

Let $(C,C')$ be good. The two links are components $Y_1,Y_2$ of $G-V(C)-V(C')$ that are joined to
both cycles, because no edge joins the cycles. The vertex $p$ lies in at most one of them. So one of
them, say $Y$, avoids $p$. Let $u$ be a vertex of $Y$ with a neighbour on $C$. The cycle $C'$ lies
in the component of $p$ in $H$, and $Y$ is connected and joined to $C'$. So $u$ lies in that
component as well. The set $Y$ is a component of $H-V(C')$ that contains $u$ and not $p$. So $C'$
separates $u$ from $p$ in $H$.

There are at most $L$ choices for $u$. So for some vertex $u$ the cycles $C'$ that separate $u$
from $p$ in $H$, as above, have total weight at least $8/L$. Take a maximal family
$C_1,\dots,C_T$ of pairwise disjoint cycles among them. Every cycle of the collection meets one of
the at most $TL$ vertices of the family. By~\eqref{eq:load},
\[
  \frac8L\;\le\;TL\cdot L\,2^{-s}.\qedhere
\]
\end{proof}

\subsection{Proof of Theorem~\ref{thm:half}}\label{sec:halfproof}

Let $D_*$ and $C_0$ be as in Lemma~\ref{lem:affine}. Let $G$ and $P$ be as in the theorem, and
suppose that $q(G,P)\ge\frac25$. We derive a contradiction when $r$ is large. All limits below are
for $r\to\infty$, and they are uniform in $G$ and $P$. Put
\[
  x=\log r,\qquad z=\log x,\qquad A=6\,e^{124},\qquad L=\lfloor Ax\rfloor,\qquad
  t=\lfloor z-2\log z\rfloor .
\]

\paragraph{Two girth scales.}
Let $D_j=2^{j}D_{*}$ and let $S_0\subseteq S_1\subseteq\cdots$ be the sets of
Theorem~\ref{thm:raise}. Let $D^{-}$ be the first $D_j$ that is at least $\sqrt x$, and let $D^{+}$
be the first $D_j$ that is at least $x/16$. For large $r$,
\[
  \sqrt x\le D^{-}<2\sqrt x,\qquad \frac x{16}\le D^{+}<\frac x8,\qquad D^{-}<D^{+}.
\]
Let $S^{-}\subseteq S^{+}$ be the corresponding sets, and put $J^{\pm}=G-(P\cup S^{\pm})$. Then
$J^{-}$ has girth greater than $D^{-}$, $J^{+}$ has girth greater than $D^{+}$, and
$J^{+}\subseteq J^{-}$. By Theorem~\ref{thm:raise}(ii),
\begin{equation}\label{eq:twosizes}
  |S^{-}|\;\le\;2^{4\sqrt x+6\log^{2}(2\sqrt x)}=r^{o(1)},\qquad
  |S^{+}|\;\le\;2^{x/4+6\log^{2}x}=r^{1/4+o(1)}.
\end{equation}
We use the small scale $D^{-}$ for the girth of the final pieces, and the large scale $D^{+}$
for a long chain.

\paragraph{Cycles at the large scale.}
We apply Section~\ref{sec:chain} with $S=S^{+}$, $J=J^{+}$, $D=D^{+}$ and
$s=\lfloor(D^{+}-1)/2\rfloor\ge x/32-1$. The hypotheses of Lemma~\ref{lem:linkshalf} hold for large
$r$: $8\le D^{+}\le L$ and $1\le t\le z\le D^{+}/4$. So every two disjoint cycles of $J^{+}$ of
length at most $L$ have at most
\[
  b\;=\;1+\frac{2L}{t+1}+5\cdot2^{t}+289\Bigl(\frac L{D^{+}}\Bigr)^{2}
  \;\le\;1+\frac{2Ax}{z-2\log z}+\frac{5x}{z^{2}}+289\,(16A)^{2}\;\le\;\frac{3Ax}{z}
\]
links. Hence
\[
  \eta=2^{b}L^{3}2^{-s}\;\le\;r^{3A/z}\,(Ax)^{3}\,2\,r^{-1/32}=r^{-1/32+o(1)}=o(1),
  \qquad L^{2}2^{-s}=o(1).
\]

The graph $J^{+}$ has $v=2r-2-|P|-|S^{+}|=2r-o(r)$ vertices, and its deficit is
$d_J\le3(|P|+|S^{+}|)=r^{1/4+o(1)}$. Let $\mu$ be the total weight of its cycles of length at most
$L$. By Lemma~\ref{lem:weight} with $\ell_0=\lceil6x\rceil$,
\[
  \Bigl(1+\tfrac32\,L^{2}\,2^{-s}\Bigr)\,\mu\;\ge\;
  \Bigl(1-\frac{L\,d_J}{2v}-3v\,r^{-3}\Bigr)\sum_{\ell}\frac1{2\ell},
\]
where the sum is over the even $\ell$ with $6x\le\ell\le L$. This sum is
$(\frac14-o(1))\ln(A/6)=31-o(1)$, and the two factors in brackets are $1+o(1)$. Thus $\mu>30$ for
large $r$.

\paragraph{The chain.}
Fix $p\in P$. For large $r$ the hypotheses of Lemma~\ref{lem:chain} hold. So there are a cycle $C$
of $J^{+}$ of length at most $L$, a vertex $u$, and pairwise disjoint cycles $C_1,\dots,C_T$ of
$J^{+}$ of length at most $L$ with
\[
  T\;\ge\;\frac{8\cdot2^{s}}{L^{3}}\;\ge\;\frac{4}{A^{3}}\cdot\frac{r^{1/32}}{x^{3}},
\]
such that each $C_i$ separates $u$ from $p$ in $G-V(C)$. Let $H$ be the component of $G-V(C)$
that contains $u$ and $p$. Each $V(C_i)$ is connected and separates $u$ from $p$, so it lies in $H$.
By Lemma~\ref{lem:nested} we may number the cycles so that $C_1<C_2<\dots<C_T$. Let $U_i$ be the
vertex set of the component of $u$ in $H-V(C_i)$.

\begin{figure}[ht]
\centering
\begin{tikzpicture}[x=1cm,y=1cm]
  % the cycle C and the vertex u
  \draw[hot] (0,0) ellipse [x radius=0.32, y radius=0.95];
  \node[note, text=warm] at (0,1.25) {$C$};
  \node[vtx] (u) at (0.85,0.25) {};
  \draw[edge] (0.31,0.25)--(u);
  \node[note, above=2pt] at (u) {$u$};
  % blocks
  \fill[accent!10] (2.4,-1.0) rectangle (5.2,1.0);
  \fill[accent!10] (5.8,-1.0) rectangle (8.6,1.0);
  \node[note, text=accent] at (3.8,-1.3) {the block $X_0$};
  \node[note, text=accent] at (7.2,-1.3) {the block $X_1$};
  % separating cycles
  \foreach \x in {2.1,5.5,8.9} \draw[cool] (\x,0) ellipse [x radius=0.26, y radius=0.95];
  \foreach \x in {2.9,3.5,4.1,4.7,6.3,6.9,7.5,8.1,9.7,10.3}
    \draw[accent, line width=0.45pt] (\x,0) ellipse [x radius=0.2, y radius=0.85];
  \node[note, text=accent] at (2.1,1.25) {$C_1$};
  \node[note, text=accent] at (5.5,1.25) {$C_{w+1}$};
  \node[note, text=accent] at (8.9,1.25) {$C_{2w+1}$};
  \node[note] at (10.95,0) {$\cdots$};
  % the marked set P
  \draw[draw=slate!70, fill=black!4, line width=0.6pt] (12.2,0) ellipse [x radius=0.55, y radius=0.45];
  \node[text=slate] at (12.2,0) {\small$P$};
  \node[vtx] (pp) at (11.92,0.05) {};
  \node[note, above=7pt] at (12.2,0.3) {$p\in P$};
  \draw[slate!60, line width=0.4pt, dash pattern=on 2pt off 2pt] (u)--(1.8,0.25) (11.3,0.05)--(pp);
\end{tikzpicture}
\caption{The chain of Lemma~\ref{lem:chain}. Each cycle $C_i$ separates $u$ from $p$ in $G-V(C)$, so
the cycles are nested. Every $w$-th cycle is the end of a block. A block has many vertices and at
most $2L$ boundary edges.}
\label{fig:chain}
\end{figure}

\paragraph{Blocks.}
Put $w=\lceil128A^{2}x\rceil+2$. For $l=0,1,\dots$ with $(l+1)w\le T-1$ let
\[
  X_l\;=\;U_{(l+1)w+1}\sm\bigl(U_{lw+1}\cup V(C_{lw+1})\bigr).
\]
There are at least $\lfloor(T-1)/w\rfloor\ge r^{1/32}/x^{5}$ such sets, for large $r$. By
Lemma~\ref{lem:blocks} they are pairwise disjoint. Each contains the $w-1$ cycles between its two
end cycles. These cycles are pairwise disjoint and each has more than $D^{+}$ vertices, so
\[
  |X_l|\;\ge\;(w-1)(D^{+}+1)\;\ge\;128A^{2}x\cdot\frac x{16}=8A^{2}x^{2}\;\ge\;8L^{2}.
\]
Every edge of $H$ that leaves $X_l$ ends on one of its two end cycles. An end cycle has at most $L$
vertices, and each of them has one edge end outside the cycle. So at most $2L$ edges of $H$ leave
$X_l$.

\paragraph{Clean blocks.}
Discard a set $X_l$ if it contains a vertex with a neighbour on $C$, or a vertex of $S^{-}$. At most
$L$ vertices outside $C$ have a neighbour on $C$, so this discards at most
$L+|S^{-}|=r^{o(1)}$ sets. For large $r$ at least two sets remain. Let $X$ be one of them.
\begin{itemize}
\item $X$ avoids $P$. The set $P$ is connected and disjoint from the cycles, and it contains $p$,
which is not in $U_i$ for any $i$. A path in $G[P]$ from $p$ to a vertex of $U_i$ would have to pass
through $C_i$ or through $C$.
\item No edge of $G$ joins $X$ to $V(C)$, and $H$ is a component of $G-V(C)$. So the edges of $G$
that leave $X$ are edges of $H$, and $|\delta X|\le2L$.
\item $X$ avoids $P\cup S^{-}$, so $G[X]$ is a subgraph of $J^{-}$, and its girth is greater than
$D^{-}$.
\end{itemize}

\paragraph{The contradiction.}
Two of the remaining sets satisfy the hypotheses of Proposition~\ref{prop:blocks}(b) with
$\Lambda=L$ and $D=D^{-}$: they are disjoint, they avoid $P$, each has at least $8L^{2}$ vertices
and at most $2L$ boundary edges, and $D^{-}\ge\sqrt x\ge C_0\log(L+2)$ for large $r$. Thus
$q(G,P)\le\frac9{25}<\frac25$. This contradiction proves Theorem~\ref{thm:half}. \qed
